% Chapter 14: Open Problems and the Next Program
% DEFINES: ch:open-problems; sec:internal-frontiers, sec:field-arrives,
%   sec:next-program, sec:closing, sec:ch14-notes; tab:directions
% REFERENCES (backward): thm:orbit-bound, rem:active-vs-marginal,
%   thm:not-approximate, rem:control-flow, thm:joint-recovery, tab:hiding,
%   sec:measured-not-negligible
% This is the book's closing chapter: a roadmap, not new results.

\chapter{Open Problems and the Next Program}
\label{ch:open-problems}

A book that builds a framework owes a closing account of what the framework
does not yet do. This chapter is that account, and it is deliberately a
roadmap rather than a results chapter: it collects the problems the
preceding parts left open, places the framework in the field that has grown
up around it, and sketches the next program of original research. The thread
that started in \Cref{sec:measured-not-negligible}, that confidentiality
should be measured rather than assumed negligible, ends here as a research
agenda.

\section{The Internal Frontiers}
\label{sec:internal-frontiers}

The sharpest open problems are the ones the framework raises about itself.
Several were flagged in place as the chapters passed them.

The active-adversary bound of \Cref{thm:orbit-bound} is loose by design: it
counts the size of the orbit but ignores the functional relations among its
elements, which constrain the latent value further than mere cardinality
does. Tightening it for specific operation sets, deriving the residual
entropy an adversary actually faces when the operations are, say, the Boolean
algebra of \Cref{ch:boolean-algebra} rather than an abstract set, is open, and
\Cref{rem:active-vs-marginal} named it so. A related question concerns noise
under composition: garbage fed through a chain of cipher maps has some
probability of landing on a valid codeword at the end, and characterizing
that noise-to-signal probability when the maps share structure, a common hash
family or overlapping codeword spaces, is unresolved; it bears directly on
the trusted machine's ability to tell a real result from an artifact. The
Boolean construction leaves its own residuals: the exact false-positive rate
of an intersection as a function of the operand sizes and bit width, and the
error rate of the approximate complement of \Cref{thm:not-approximate} as a
function of the fill factor, are both open in closed form. And the cost of
oblivious control flow, the choice \Cref{rem:control-flow} forced between
leaking a test and evaluating both branches, has a decision space, which cuts
to make, how the ciphered region propagates, that the program-realization
work only begins to map.

These are not gaps that threaten the framework; they are the questions a
working theory generates. Each is a concrete, well-posed problem with a clear
success criterion, which is the mark of a frontier worth working.

\section{The Field Arrives}
\label{sec:field-arrives}

When this program began, its central bet, that a measurable, tunable leakage
sitting between plaintext and the heavy machinery of oblivious RAM or fully
homomorphic encryption is a legitimate design point, was contrarian. The
field has since moved toward it. The searchable-encryption literature, long
dominated by leakage-abuse attacks, pivoted around 2024 toward
quantify-then-mitigate, with systems that name and bound leakage rather than
hoping it is negligible; the ``tunable measurable leakage'' middle ground
crystallized into a recognized design space in the same window. The
quantitative-information-flow community matured a family of operational
leakage measures, maximal leakage and its relatives, whose live threads,
composition and correlated secrets, are exactly the questions the
compositional scale of \Cref{ch:compositional-scale} addresses. A subfield on
adversarially-robust probabilistic data structures consolidated, and the
cipher map is a native citizen of it with a privacy notion the subfield has
not yet named. And in private machine-learning inference, the idea that
approximation itself can be privacy, the dual reading of the error budget
$\eta$, is being rediscovered independently.

The lesson is not that the program was right in every claim. The honest
assessment, recorded in the program's own survey, is that its strongest
theorems are sound but that several papers reached for a broader qualitative
``first'' that an adjacent vocabulary, sampling theory, distribution
matching, wire-tap coding, had anticipated. The lesson is rather that the
program's taste in \emph{problems} was good, and is now externally
validated: being early to a space the field later judged important is the
position to be in.

\section{The Next Program}
\label{sec:next-program}

What the framework should do next follows from one observation about where it
is uniquely strong. Every neighboring system in the tunable-leakage space
pays for its hiding \emph{per query}, in bandwidth, rounds, or injected
noise. The codec construction of \Cref{ch:gf2-retrieval} does not: its
frequency-hiding is structural, built once and paid for in a rank condition,
at zero marginal cost per access (\Cref{tab:hiding}). That single property,
structural frequency-hiding at zero per-query cost in a static data
structure, is the program's clearest contribution to the field, and most of
the salient directions extend it (\Cref{tab:directions}).

\begin{table}[t]
\centering
\caption{Salient directions for the next program.}
\label{tab:directions}
\begin{tabular}{@{}p{0.40\textwidth}p{0.52\textwidth}@{}}
\toprule
Direction & The gap it fills \\
\midrule
The non-member output law as a privacy axiom for probabilistic data structures
  & PDS correctness models do not treat the false-positive answer law as a controllable, frequency-independent design object \\
The zero-per-query point on the tunable-leakage frontier
  & every system in the 2024 tunable-leakage space pays per query; this one does not \\
Re-expressing the leakage knob in maximal- / alpha-leakage currency
  & operational-leakage composition with correlated secrets is the active QIF thread \\
Adaptive multiplicity retuning under distributional drift
  & nobody owns maintaining a quantified leakage bound while the query distribution drifts \\
A codec-controlled membership layer for private retrieval and vector search
  & semantic-search systems pay oblivious-RAM or homomorphic costs; embedding inversion is a recognized top risk \\
\bottomrule
\end{tabular}
\end{table}

The directions are not equally ripe. The first two are framing-level: they
convert results already proved in this book into conversations the field is
actively having, and cost a table and a careful positioning rather than a new
theorem. The third buys insurance against the simulation-based reviewer, by
speaking the leakage measure in the currency that community has standardized.
The fourth, maintaining a leakage bound as the distribution drifts, is the
natural next paper and the cleanest open lane, since no one yet owns
quantified-leakage-under-drift. The fifth is the largest audience and the
largest lift: a private-retrieval substrate for the machine-learning systems
that have made encrypted semantic search urgent.

\section{A Closing Word}
\label{sec:closing}

The argument of this book has been a single one, developed in six movements.
A cipher map is a total function on bit strings that an untrusted machine
evaluates blindly; its four properties make it measurable rather than
magical; its algebra and types say what can and cannot be computed through
the trapdoor; its confidentiality is a quantity, governed by $\delta$ at the
margin and by something $\delta$ cannot touch under composition; it is
realized by concrete constructions, one of which hides frequency for free;
and it opens onto a research program the field has begun to converge on
independently.

The framework's privacy does not come from hiding access patterns, and it is
not oblivious RAM, homomorphic encryption, or a simulation-based security
game. It comes from a one-way hash and a uniform representation, and its
guarantees are numbers a designer can read and trade. That is a narrower
promise than cryptography's strongest tools make, and a more honest one for
the systems that cannot afford them: a service that cannot read its queries
or its answers, returns the right ones anyway, and leaks an amount you can
write down. The next program is to make that promise sharper, to speak it in
the field's own measures, and to build it into the systems that need it. The
foundation is here.

\section{Notes and Provenance}
\label{sec:ch14-notes}

This chapter is a roadmap, not a results chapter, and claims no new theorems.
The internal open problems are collected from the points where the earlier
chapters flagged them and from the open-problems sections of the construction
papers. The field landscape and the ranked directions are drawn from the
program's 2026 ecosystem survey, which mapped the searchable-encryption,
quantitative-information-flow, adversarial-data-structure, and
private-inference literatures against the program and assessed its standing
frankly; the external systems and threads named here are pointed to by name
and venue rather than cited formally, pending the citation pass that will
fold them into the bibliography. The unifying observation, structural
frequency-hiding at zero per-query cost as the program's distinctive
contribution, is that survey's, and it is the thread the next program should
pull.
